% Source: https://artofproblemsolving.com/community/c3h3267911_2023_lmt_fall_guts_round_p1p15_lexington_mathematical_tournament
% Stable identity: aops-topic-3267911-q15
Only critical radii matter; touching a square at a point or along a line does not count. Arbitrarily small positive disks centered at a grid vertex cover four squares, so $a_1=a_2=a_3=a_4=0$. A disk of radius at most $\dfrac{1}{2}$ cannot cover three columns or three rows, so it covers at most four squares. Centering at the midpoint of a grid edge and increasing the radius just beyond $\dfrac{1}{2}$ covers six. Thus $a_5=a_6=\dfrac{1}{2}$. For radii below 1, at most a $3\times3$ block can be covered. To cover seven squares, at least two of its four corner squares must be reached. Opposite corners require radius greater than $\dfrac{\sqrt2}{2}$. If exactly two adjacent corners are reached, all five noncorner squares must also be reached. Put the central cell at $[0,1]^2$ and the disk center at $(x,y)$, with the reached corners below it. The necessary radius is greater than both $1-y$ and $\sqrt{\dfrac{1}{4}+y^2}$, because the farther of $(0,0),(1,0)$ is at least that far away. Minimizing their maximum gives $y=\dfrac{3}{8}$ and radius $\dfrac{5}{8}$. At center $(\dfrac{1}{2},\dfrac{3}{8})$, a radius just above $\dfrac{5}{8}$ indeed reaches seven squares. Hence $a_7=\dfrac{5}{8}$. Reaching eight squares requires at least three corners, including an opposite pair, so radius must exceed $\dfrac{\sqrt2}{2}$; a disk centered at a cell center then reaches all nine. Thus $a_8=a_9=\dfrac{\sqrt2}{2}$. Radius at most 1 can cover at most three columns and three rows, hence at most nine squares. Just above 1, a disk centered at a grid vertex reaches twelve, so $a_{10}=a_{11}=a_{12}=1$. Finally $a_{13}>1$: even a radius of $\dfrac{11}{10}$ cannot reach thirteen squares. Its bounding box meets at most four rows and four columns. Thirteen squares would require all four of each and a corner square of that $4\times4$ block. After translating its central grid vertex to the origin, the disk center must have both coordinates of magnitude below $\dfrac{1}{10}$. Every corner square is then at distance at least $\dfrac{9\sqrt2}{10}>\dfrac{11}{10}$, impossible. The changes therefore occur at $n=4,6,7,9,12$, whose sum is $\boxed{38}$.
\begin{asy}
import geometry;
usepackage("amsmath");
size(300);
real[] rr={.5,.625,sqrt(2)/2,1};pair[] cc={(.5,0),(.5,.375),(.5,.5),(0,0)};string[] tt={"$a_5=a_6=\dfrac{1}{2}$","$a_7=\dfrac{5}{8}$","$a_8=a_9=\dfrac{\sqrt2}{2}$","$a_{10}=a_{11}=a_{12}=1$"};for(int k=0;k<4;++k){pair off=(4*k,0);for(int i=-1;i<=2;++i){draw(off+(i,-1)--off+(i,2),lightgray);draw(off+(-1,i)--off+(2,i),lightgray);}draw(circle(off+cc[k],rr[k]),blue);dot(off+cc[k]);label(tt[k],off+(.5,-1.5));}
\end{asy}
